% !TeX program = xelatex
% !TeX encoding = UTF-8
\documentclass{MathorCupmodeling}
\bianhao{XXXXXX}
\tihao{（A/B/C/D）}
\timu{论文题目}
\keyword{content；content；content}
\begin{document}
	\begin{abstract}
		[中文摘要内容：600--800字，写满且收在一页；每个子问题独立一段，包含具名方法和真实数值结果，并用 \textbf{} 各突出1--3处关键方法/结果]
	\end{abstract}
	\newpage
	% 中文多问题竞赛论文的默认前置结构：重述说明“做什么”，分析按题号说明“为何这样建模”。
	% 写作时可按实际问题数增删主体章节，但不要删除逐问问题分析的语义结构。
	\input{sections/1_restatement}
	\input{sections/2_analysis}
	\input{sections/3_assumptions}
	\input{sections/4_symbols}
	\input{sections/5_problem1}
	\input{sections/6_problem2}
	\input{sections/7_problem3}
	\input{sections/8_sensitivity}
	\input{sections/9_evaluation}

	\phantomsection
	\addcontentsline{toc}{section}{参考文献}
	\begin{thebibliography}{99}
	\bibitem{label}content
	\end{thebibliography}

	\newpage
	\appendix
	\ctexset{section={
		format={\zihao{-4}\heiti\raggedright}
	}}
	\begin{center}
		\heiti\zihao{4} 附\hspace{1pc}录
	\end{center}
	\input{sections/A_code}
\end{document}
